% Source: 06 - Wed Oct 7 2026.pdf, page 4. Content remains in source order.
\sourcepage{06}{4}
We introduce two extra problems on graphs:

\begin{notesdefinition}
\textbf{DEF:} Given $G=(V,E)$ undirected, a vertex cover of size $k$ is a subset of nodes $V'\subseteq V$, $|V'|=k$ such that each edge in $E$ has at least one endpoint in $V'$.
\[
\begin{gathered}
V'\subseteq V,\ |V'|=k:\ \forall u\neq v\in V:\\
\{u,v\}\in E\Rightarrow(u\in V')\lor(v\in V').
\end{gathered}
\]
\end{notesdefinition}

\textbf{EXAMPLE:}
\sourcefigure[0.7]{assets/06/page-004-vertex-cover.png}
We have:

\textbf{VERTEX COVER (VC)}

\begin{notesproblem}
\text{I:} & \langle G=(V,E),k\rangle\text{, }G\text{ undirected graph, }1\leq k\leq|V|\text{.}\\
\text{Q:} & \text{Does }G\text{ contain a vertex cover of size }k\text{?}
\end{notesproblem}

The optimization version of this problem requires determining the vertex cover of min size.

A vertex cover identifies a (possibly small) set of nodes from where all edges can be controlled. Crucial in networking, distributed computing, etc.\ldots
